\chapter{Revised Plan and Remaining Work}
\label{app:revised-plan}

The revised July 27--October 27, 2026 plan requires V2 retesting, loss and temperature diagnosis, and V3 development \cite{projectOctoberPlan}. Section~\ref{sec:res-revised-plan} gives the current status.

\section{Plan Documents and Power Requirements}

The earlier plan included a \SI{75}{\watt} continuous-validation milestone \cite{projectEarlierPlan}.

The revised plan sets a \SI{65}{\watt} continuous minimum and keeps the \SI{100}{\watt} continuous target. The reported point at approximately \SI{85}{\watt} for V2 still needs to be repeated or explained. V2 includes the integrated MCU and the ported control firmware; the modified V1.0 data are from the earlier assembly. The board, controller, and settings must be recorded with the measurements when the V2 test is repeated.

Section~\ref{sec:disc-plan-gap} explains the \SI{15}{\watt} and \SI{20.11}{\watt} output shortfalls and the \SI{24.53}{\watt} screening input--output difference. The sustained-output limit needs to be identified in the next tests.

\section{Thesis Sections for the Revised Plan}

\begingroup
\small
\begin{longtable}{@{}L{0.27\linewidth}L{0.65\linewidth}@{}}
\caption[Revised plan and thesis sections.]{Revised plan and corresponding thesis sections.}
\label{tab:plan-thesis-map}\\
\toprule Planned work & Thesis section\\\midrule
\endfirsthead
\toprule Planned work & Thesis section\\\midrule
\endhead
Objectives and test criteria & Chapter~\ref{ch:introduction} sets the revised scope. Section~\ref{sec:revised-plan-criteria} defines AP-01--AP-08, including power levels, tests, and design documents.\\
Repeat V2 and diagnose the gap & Section~\ref{sec:method-plan-validation} describes how to repeat the baseline, measure input and output power, increase the load in stages, and change one parameter at a time. Section~\ref{sec:disc-plan-gap} discusses possible causes of the power limit.\\
V2 circuit and software & Sections~\ref{sec:hw-v2} and \ref{sec:ctrl-pwr2-port} document the V2 circuit and MCU firmware. Appendices~\ref{app:firmware} and \ref{app:software-tools} explain the software; the complete source follows in Appendices~\ref{app:full-embedded-source} and \ref{app:full-supporting-source}.\\
V3 edge-card implementation & Section~\ref{sec:hw-v3-plan} describes the reduced edge-card layout, removal of the USB-PD section, and the replacement magnetics and hardware checks still needed. Section~\ref{sec:res-v3-design} reports the design outcome.\\
PDO/PPS, thermal, and stability tests & Section~\ref{sec:method-plan-validation} identifies the required protocol, dynamic, protection, and defined-duration tests. Section~\ref{sec:res-revised-plan} summarizes the available test results.\\
Power density and fallback & Section~\ref{sec:method-density} defines a common assembly-volume and continuous-power basis. Section~\ref{sec:res-revised-plan} reports the V3 design-volume estimate; a complete-system power-density comparison and readiness decision remain unfinished.\\
\bottomrule
\end{longtable}
\endgroup

\section{Planned Dates and Required Work}

Table~\ref{tab:plan-milestones} lists the revised plan's technical milestones and submission dates. Section~\ref{sec:res-revised-plan} summarizes the completed designs and available test results.

\begingroup
\small
\begin{longtable}{@{}L{0.17\linewidth}L{0.75\linewidth}@{}}
\caption[Revised-plan milestones in 2026.]{Scheduled milestones in the revised July--October 2026 plan. All dates are in 2026.}
\label{tab:plan-milestones}\\
\toprule Planned date & Planned work\\\midrule
\endfirsthead
\toprule Planned date & Planned work\\\midrule
\endhead
August 7 & Updated plan and test procedure: operating and stop limits, instruments, test matrix, raw-file locations, and common V2/V3 volume convention.\\
August 21 & V2 hardware and bench checked; the approximately \SI{85}{\watt} point repeated or its difference explained.\\
September 4 & V2 operating-point test set completed, with raw files and configuration records.\\
September 11 & Power-gap analysis, uncertainty, loss/stress assessment, and changes affecting the V3 fabrication files completed.\\
September 16 & V3 schematic/layout, design-rule, BOM, and manufacturing reviews completed; the board released with its size and density definition. Fabrication was scheduled for September 17.\\
September 25 & Adjusted V2 retested, with measured effects of the selected changes.\\
September 30 & Target V3 board arrival; October 2 was the latest planned arrival. Assembly, inspection, and gradual bring-up were scheduled before the readiness decision.\\
October 7 & V3 readiness decision: proceed with matched final V3 tests, or use completed, identified V2 tests and explain the remaining V3 work.\\
October 14 & Final tests completed: sustained \SI{100}{\watt} under stated conditions, or the highest repeatable supported power and remaining limit, together with the agreed electrical, thermal, stress, and density results.\\
October 16 & Complete thesis draft with methods, V2 findings, V3 design/status, results, figures, limitations, and conclusions.\\
October 19--26 & Supervisor corrections, technical-claim and reference checks, final compilation, and visual PDF review. Final submission files were scheduled to be ready on October 26.\\
October 27 & Thesis submission.\\
\bottomrule
\end{longtable}
\endgroup

After the scheduled September 16 release, the plan limits changes to tuning, firmware, compatible parts, testing, and controlled rework unless a safety problem requires new fabrication files. This keeps the final test period focused on the released design. Further PCB ideas belong in the next revision.

The matched numerical measurements are from the earlier prototype. V2 adds the integrated circuit, software, and later tank and thermal observations. V3 adds a slimmed layout with secondary rectification retained and the USB-PD section removed. A replacement magnetic set was not installed, and a matched complete-system V2--V3 power-density comparison remains unfinished.